\documentclass[11pt]{article}
\usepackage{amsmath,amssymb,amsthm}
\usepackage[margin=1in]{geometry}
\usepackage{url}
% Allow line breaks in paths/URLs at many characters (no xurl dependency).
\makeatletter
\g@addto@macro{\UrlBreaks}{\do\/\do\-\do\.\do\_\do\~\do\=\do\&\do\?\do\%}
\def\UrlFont{\ttfamily}
\makeatother
\usepackage[hidelinks,breaklinks]{hyperref}
\hypersetup{
  pdftitle={Poisson and Weyl lifts of an explicit three-dimensional Keller map and nonunique self-adjoint realizations of a dual transport operator},
  pdfauthor={Daniel Eric Fredriksen},
  pdfsubject={Deficiency indices and self-adjoint extensions of a dual transport operator},
  pdfkeywords={Jacobian conjecture, deficiency indices, self-adjoint extensions, Weyl algebra}
}
% Breakable monospaced path/filename (wraps like \url).
\newcommand{\filepath}[1]{\path{#1}}
\newtheorem{theorem}{Theorem}
\newtheorem{proposition}{Proposition}
\newtheorem{lemma}{Lemma}
\newtheorem{corollary}{Corollary}
\newtheorem{remark}{Remark}
\title{Poisson and Weyl lifts of an explicit three-dimensional Keller map\\
and nonunique self-adjoint realizations of a dual transport operator}
\author{Daniel Eric Fredriksen\\
Quantyra Inc.\\
\url{https://github.com/Quantyra/jacobian-weyl-quantum-phase-space}}
\date{29 July 2026}

\begin{document}
\maketitle

\begin{abstract}
We study the three-dimensional Keller map \(F\) announced by Alp\"oge, whose announcement credits Fable for work leading to the map: \(\det DF\equiv -2\) and a three-point collision. Theorem~A restates those finite identities. Theorems~B--C record the standard cotangent/Piola/Weyl package for this seed; the analytic contribution is Theorems~D--F. The cotangent lift \(\Phi(q,p)=(F(q),J^{-T}p)\) is a polynomial Poisson map and non-injective. Dual fields \(X_j=\sum_k B_{jk}\partial_{q_k}\) commute and are divergence-free (Piola), so \(H_j=-iX_j\) are symmetric on \(C_c^\infty(\mathbb{R}^3)\). We define the induced polynomial Weyl endomorphism \(\psi\), without claiming that \(\psi\) is noninvertible or not an automorphism. The field \(X_1\) is incomplete via an explicit integral curve. In flow-box coordinates \((a,s,c)=F(q)\) one has \(X_1=\partial_s\) and \(\mathrm{d}q=\tfrac12\,\mathrm{d}a\,\mathrm{d}s\,\mathrm{d}c\). An open forward escape wall yields whole-orbit vectors in \(\ker(H^*+i)\) (transverse cutoffs only; no interior \(s\)-indicators); a measure-preserving sign involution transports them isometrically to \(\ker(H^*-i)\). Thus the Hilbert deficiency indices are
\[
(n_+,n_-)=(\aleph_0,\aleph_0),
\]
equivalently \((\infty,\infty)\) in conventional deficiency-index notation, for \(H=-iX_1\). A continuum-sized family of distinct self-adjoint extensions exists; the core algebraic CCR/Poisson data alone do not distinguish among them. The operator-theory mechanism is classical; the contribution is its explicit realization for this dual field. No gates, channels, computational advantage, or seed discovery claim. Paper-package candidate \textbf{v0.3.9-referee-revision} (proposed, untagged).
\end{abstract}

\noindent\textbf{MSC 2020:} 47B25, 81S05, 14R15, 37C10, 53D17.\\
\textbf{Paper package:} v0.3.9-referee-revision (proposed, untagged); concept DOI\\
\url{https://doi.org/10.5281/zenodo.21474351}
(project-level; no new version DOI is asserted before release ingest).\\
\textbf{Categories (suggested):} math.SP (primary); math.FA, math-ph (secondary).

\paragraph{Lean coverage boundary.}
\begingroup
\emergencystretch=2em
The companion Lean artifact proves the bounded canonical-core declaration
\texttt{ExoticCCR.theoremE}: the minimal transport core for \(X_1\) is not
essentially self-adjoint via a constructed weak \(-i\) witness. The historical
release tag \texttt{v0.1.8-theorem-e} targets commit \texttt{be4f330};
\texttt{a6bb091} is the reviewed Theorem~E source anchor within that release
history, not the tag target. The synchronized Lean
publication freeze is
{\small\texttt{fbcdd0345d2f2540cd537204be2178ae07e18a5e}}; it is untagged and
unreleased. The earlier commit
\texttt{ff50f4a2a312591c2e5b26e71eb390ade9164b34} is retained only as the
historical Theorem~F source root. At the synchronized freeze the artifact proves
the bounded Theorem~F declaration \texttt{ExoticCCR.theoremF}: both adjoint
eigenspaces $\ker(H^\dagger \mp i)$ of the canonical minimal core are not
finite-dimensional and the bounded reading of
$(n_+,n_-)=(\infty,\infty)$. It also proves \(\operatorname{aleph0}\) lower
bounds for the corresponding \texttt{Module.rank}-valued algebraic/Hamel
ranks. The
Lean route to the $+i$ side is a measure-preserving sign-involution transport
($\sigma(q_0,q_1,q_2)=(-q_0,-q_1,q_2)$, $X_1\circ\sigma=-\sigma_*X_1$) of the
forward minus-i families, not this paper's backward-wall construction; the
same SHA separately defines \texttt{hilbertDeficiencyIndex} as the cardinality
of a chosen Hilbert basis for each closed adjoint eigenspace and proves the
exact-index application theorem
\texttt{hilbertDeficiencyIndex}\allowbreak\texttt{\_X1\_eq\_aleph0}; thus the exact Hilbert-space
indices are $\operatorname{aleph0}$ on both sides. This is distinct from
\texttt{hamelDeficiencyRank}, which is \texttt{Module.rank}
(algebraic/Hamel rank) and has lower bounds only. For the specific canonical
minimal operator \texttt{H\_X1\_min}, the same SHA proves a bijective von
Neumann classification of all
\texttt{SelfAdjoint}\allowbreak\texttt{Extension H\_X1\_min}
witnesses (self-adjoint \texttt{LinearPMap}s extending
\texttt{H\_X1\_min}) by complex-linear isometric equivalences from the $+i$
adjoint eigenspace to the $-i$ adjoint eigenspace. The declarations
\texttt{theoremF}\allowbreak\texttt{UnitPhase}\allowbreak\texttt{Extension} and
\texttt{theoremF}\allowbreak\texttt{UnitPhase}\allowbreak\texttt{Extension\_injective}
prove that unitary complex
phases inject into distinct such witnesses. The further conclusion that this
is a continuum-sized lower family uses the classical cardinality of the
complex unit circle; that cardinal identification is not part of these Lean
declarations~\cite{FredriksenLeanF}.

The paper's independent backward-wall proposition, the exact cardinality of
the full extension type, preferred extension, arbitrary-operator
classification, and algebraic/physical selection claims remain outside the
Lean corollary.
Formal symmetry of the specific canonical minimal core
\texttt{H\_X1\_min} on its $C_c^\infty$ test-function domain is separately
Lean-formalized at Stage~A SHA \texttt{bc679df} in
\texttt{TheoremF}\allowbreak\texttt{SymmetricCore}, with the corresponding
graph-space foundations in
\texttt{TheoremF}\allowbreak\texttt{VonNeumannGraph}. This is not an arbitrary-operator or
arbitrary-transport-realization symmetry theorem.
\par
\endgroup

\paragraph{Errata (prior releases).}
v0.2.2 pair \((0,\infty)\) for \(H\) withdrawn. v0.3.0 interior \(s\)-indicator deficiency proof withdrawn (\(\notin\operatorname{Dom}(H^*)\)). This submission is \textbf{\(H=-iX_1\) only}: we do \emph{not} claim exact deficiency pairs for \(H_0,H_2\), nor a theorem-level dual-flow strong CCR no-go (companion notes only).

\section{Introduction}
On 19--20 July 2026, Levent Alp\"oge announced an explicit polynomial map \(F:\mathbb{C}^3\to\mathbb{C}^3\) with \(\det DF\equiv-2\) and a three-point collision. The announcement credits Akhil Mathew with prompting the question and Claude Fable with work leading to the map~\cite{Alpoge2026,AFPJacobian}. \textbf{We do not claim discovery of \(F\).} The finite seed identities are \emph{verified} here by dual CAS and by a companion Lean 4 formalization~\cite{FredriksenLean,FredriksenSoft}; an independent Isabelle/HOL development verifies the determinant and collision and records the same attribution~\cite{AFPJacobian}. Our contribution is the dual-lift deficiency analysis (Theorems~D--F), not the Jacobian-conjecture refutation itself.

We ask whether the standard Poisson/CCR lifts of \(F\) select a unique self-adjoint dual transport generator on \(C_c^\infty(\mathbb{R}^3)\). For \(H=-iX_1\) the answer is negative: deficiency indices \((\infty,\infty)\).

\paragraph{Scope.} Theorems A--F below concern the algebra of the lifts and the analytic deficiency theory of \(H=-iX_1\). Remarks on other dual fields or alternative realizations are deferred to companion validation notes and are not claimed as theorems here.

\subsection*{Algebraic and geometric context}
The cotangent, Piola, and Weyl constructions in Theorems~B--C are standard
once a Keller map is given; they are included to fix conventions and connect
the seed to the analytic problem, not claimed as a new construction. Their
Weyl-algebra setting is related to the Dixmier conjecture; for classical
Jacobian-conjecture background see~\cite{BassCW}. The classical
same-rank implication is \(\mathrm{DC}_n\Rightarrow\mathrm{JC}_n\), whereas
Tsuchimoto and, independently, Belov--Kanel--Kontsevich proved the
dimension-doubling converse
\(\mathrm{JC}_{2n}\Rightarrow\mathrm{DC}_n\); together these give stable
equivalence~\cite{Tsuchimoto2005,BKK2007,AdjamagboEssen2007}. In particular,
one must not replace the latter by a dimension-preserving
\(\mathrm{JC}_n\Rightarrow\mathrm{DC}_n\). We do not derive or claim here a
new rank-three Dixmier counterexample, nor do we claim that the explicitly
defined \(\psi\) is noninvertible.

The July 2026 follow-on literature is developing rapidly. The Ulam preprint
analyzes the seed's fibers and nonproperness~\cite{UlamJacobian}; Shaska studies
its mixed-sign grading, quotient geometry, discriminant, and nonproperness
mechanism~\cite{Shaska2026}; Speyer records an early public algebraic
discussion~\cite{SpeyerJC2026}; and the Isabelle/AFP entry supplies an independent
formal check of the finite seed identities~\cite{AFPJacobian}. These recent
sources are cited as current preprints or formal artifacts, not as completed
peer review. More generally, nonproperness sets of generically finite
polynomial maps have an established literature~\cite{JelonekLason2018}. The
escape wall below is an operator-theoretic use of explicit behavior at
infinity; it does not purport to classify the full nonproperness set.

\section{Setup}
Let \(q=(q_0,q_1,q_2)=(x,y,z)\) and
\begin{align*}
F_0&=(1+xy)^3 z + y^2(1+xy)(4+3xy),\\
F_1&=y + 3x(1+xy)^2 z + 3xy^2(4+3xy),\\
F_2&=2x - 3x^2 y - x^3 z.
\end{align*}
Write \(J=DF\) and \(B:=J^{-T}\). Dual fields:
\[
X_j:=\sum_{k=0}^{2} B_{jk}(q)\,\partial_{q_k},\qquad j=0,1,2.
\]
\paragraph{Normalizations.}
We use two compatible packages. (i)~\emph{Abstract Weyl algebra} \(\mathcal{W}\): generators \(q_k,p_\ell\) with \([q_i,p_j]=\delta_{ij}\) (Dixmier convention). (ii)~\emph{Schr\"odinger representation} on \(L^2(\mathbb{R}^3)\): multiplication by \(q_k\) and \(p_k^{\mathrm{Sch}}=-i\partial_{q_k}\), so \([q_i,p_j^{\mathrm{Sch}}]=i\delta_{ij}\). The operators \(H_j:=-iX_j\) live in package~(ii). When \(\operatorname{div} X_j=0\), each \(H_j\) is symmetric on \(C_c^\infty(\mathbb{R}^3)\). Our complex inner products are conjugate-linear in the first argument and linear in the second. Collision data (char \(\neq 2\)):
\begin{equation}
F\bigl(0,0,-\tfrac14\bigr)
=F\bigl(1,-\tfrac32,\tfrac{13}2\bigr)
=F\bigl(-1,\tfrac32,\tfrac{13}2\bigr)
=\bigl(-\tfrac14,0,0\bigr).
\label{eq:collision}
\end{equation}

\section{Theorem A --- seed identities}
\begin{theorem}[A]
(1) \(\det DF\equiv -2\).
(2) The collision identities \eqref{eq:collision} hold over any field of characteristic not \(2\).
(3) Hence \(F\) is not injective.
\end{theorem}
\begin{proof}
Expand \(\det J\) as a polynomial identity, or evaluate the formal determinant in a computer algebra system; formalized in Lean~4 in the companion repository
(\filepath{ExoticCCR/AnchorF.lean}: \texttt{jacobianDet\_F}, \texttt{evalMap\_F\_p*})~\cite{FredriksenLean}.
Non-injectivity is immediate from \eqref{eq:collision}.
\end{proof}

\begin{remark}[Complex vs real]
The seed is stated algebraically over \(\mathbb{C}^3\), but all three collision
points in \eqref{eq:collision} are rational and hence real. Thus the same
polynomial formulas restrict to a noninjective local diffeomorphism of
\(\mathbb{R}^3\). We treat \(F\) as an externally announced
seed~\cite{Alpoge2026,UlamJacobian,AFPJacobian}; Theorem~A is a restatement
with CAS/Lean re-verification, not a priority claim. Theorems~D--F concern this
real restriction and the associated \(L^2(\mathbb{R}^3)\) dual-field package.
\end{remark}

\section{Theorem B --- Poisson cotangent lift}
\begin{theorem}[B]
Let \(\Phi(q,p)=(F(q),B(q)p)\) with \(B=J^{-T}\). Then:
\begin{enumerate}
\item \(B\) is polynomial and \(JB^{T}=I\);
\item writing \(Q_i=F_i(q)\) and \(P_j=\sum_k B_{jk}(q)p_k\),
\[
\{Q_i,Q_j\}=0,\quad
\{P_i,P_j\}=0,\quad
\{Q_i,P_j\}=\delta_{ij}
\]
for the canonical Poisson structure on \(T^*\mathbb{R}^3\);
\item \(\Phi\) is not injective (lift the three collision sources with matching \(p\)).
\end{enumerate}
\end{theorem}
\begin{proof}
(1) follows from \(B=J^{-T}\) and \(\det J\equiv-2\neq 0\); polynomiality of \(B\) is the classical adjugate formula \(B=-(\operatorname{adj} J)^{T}/2\).
(2) The identities \(\{Q_i,P_j\}=(JB^{T})_{ij}=\delta_{ij}\) are immediate; vanishing of \(\{Q_i,Q_j\}\) is clear; vanishing of \(\{P_i,P_j\}\) is the standard cotangent-lift calculation for \(B=J^{-T}\) (equivalently dual-CAS generator checks in the companion release~\cite{FredriksenSoft}).
(3) If \(F(q^{(a)})=Q_\star\) and \(p^{(a)}=J(q^{(a)})^{T}P_\star\), then \(\Phi(q^{(a)},p^{(a)})=(Q_\star,P_\star)\) for each collision source in \eqref{eq:collision}.
\end{proof}

\section{Theorem C --- dual fields, Piola, Weyl endomorphism}
\begin{theorem}[C]
\begin{enumerate}
\item \(X_j(F_i)=\delta_{ij}\), \([X_i,X_j]=0\), and \(\operatorname{div} X_j=0\).
\item Each \(H_j=-iX_j\) is symmetric on \(C_c^\infty(\mathbb{R}^3)\), and \([M_{F_i},H_j]=i\delta_{ij}\) on that domain.
\item Let \(\mathcal{W}\) be the polynomial Weyl algebra over \(\mathbb{C}\) generated by \(q_k,p_\ell\) with \([q_i,p_j]=\delta_{ij}\) (abstract package) and all other generators commuting. The assignment
\[
\psi(q_i)=F_i(q),\qquad
\psi(p_j)=\sum_k B_{jk}(q)\,p_k
\]
extends uniquely to a unital algebra endomorphism \(\psi:\mathcal{W}\to\mathcal{W}\) preserving the Weyl relations on generators (ordering: coefficients \(B_{jk}(q)\) to the left of \(p_k\)).
\end{enumerate}
\end{theorem}
\begin{proof}
(1) \(X_j(F_i)=(JB^{T})_{ij}=\delta_{ij}\). Then \([X_i,X_j](F_k)=0\) for all \(k\); since \(\mathrm{d}F\) is a coframe, \([X_i,X_j]=0\). Piola: \(\operatorname{Div}\operatorname{Cof}(DF)=0\); with \(\det DF\equiv-2\), \(\operatorname{Div} B=0\), i.e.\ \(\operatorname{div} X_j=0\).
(2) Divergence-free smooth vector fields give symmetric \(-iX\) on \(C_c^\infty\)~\cite{RS2}; the multiplication commutator is \(iX_j(F_i)\) in the Schr\"odinger package.
(3) By the universal property of \(\mathcal{W}\)~\cite{DixmierWeyl} it is enough to check the images of generators satisfy the same relations. Directly (no Poisson-to-commutator dictionary): with left-ordered coefficients,
\begin{align*}
\bigl[\psi(q_i),\psi(p_j)\bigr]
&=X_j(F_i)=\delta_{ij},\\
\bigl[\psi(p_i),\psi(p_j)\bigr]
&=\text{the first-order operator of }[X_i,X_j]=0.
\end{align*}
(The identities are algebraic in the Weyl algebra and match the dual-CAS generator checks in~\cite{FredriksenSoft}.)
\end{proof}

\paragraph{Lean coverage for Theorem~C(2).}
The all-\(j\) analytic symmetry and commutator assertion above
is the classical paper-level consequence of the Lean-verified polynomial
divergence and dual-field identities. The pinned operator-level Lean theorem
\filepath{H_X1_min_isFormalAdjoint_H_X1_min} formalizes the specific
\(j=1\) compactly supported core used in Theorems~E--F; the freeze does not
claim a named all-\(j\) operator theorem.

\begin{remark}
\(\psi\) is \emph{not} claimed to extend to a unitary conjugation on \(L^2\), a \(C^*\) automorphism, or a quantum channel.
\end{remark}

\section{Theorem D --- incompleteness of \texorpdfstring{\(X_1\)}{X1}}
\begin{theorem}[D]
The dual field \(X_1\) is incomplete on \(\mathbb{R}^3\).
\end{theorem}
\begin{proof}
For \(t\in[0,\tfrac12)\), with continuous extension \(\gamma(0)=(1,0,0)\), set
\begin{align*}
q_0(t)&=\frac{-2t-\sqrt{1-2t}+1}{t(2t-1)},\\
q_1(t)&=t,\\
q_2(t)&=t^2\bigl(2t-3\sqrt{1-2t}-1\bigr)
\end{align*}
on \((0,\tfrac12)\), and \(\gamma(t)=(q_0(t),q_1(t),q_2(t))\). Then \(\gamma\) is \(C^\infty\) on \((0,\tfrac12)\) and extends \(C^1\) to \(t=0\). Computer algebra gives the polynomial/algebraic identity \(F(\gamma(t))=(0,t,2)\) on \((0,\tfrac12)\)~\cite{FredriksenSoft}. Differentiating yields \(DF_{\gamma(t)}\gamma'(t)=e_1\). Since \(DF\cdot X_1=e_1\) uniquely, \(\gamma'=X_1\circ\gamma\).

As \(t=\tfrac12-\varepsilon\) with \(\varepsilon\downarrow 0\),
\[
q_0(t)=\sqrt{2/\varepsilon}+O(1)\to+\infty,
\]
so \(\gamma(t)\) leaves every compact of \(\mathbb{R}^3\) in finite time. Hence \(X_1\) is incomplete.
\end{proof}

\section{Flow-box coordinates}
Along the \(X_1\) flow (where defined), \(X_1 F_0=X_1 F_2=0\) and \(X_1 F_1=1\). On any open set carrying a smooth local inverse of \(F\) (locally everywhere: \(\det DF\equiv-2\neq 0\)), the coordinates \((a,s,c)=F(q)\) satisfy
\begingroup
\renewcommand*{\theHequation}{flowbox}
\begin{equation}
\mathrm{d}q
=\frac{1}{|\det DF|}\,\mathrm{d}a\,\mathrm{d}s\,\mathrm{d}c
=\tfrac12\,\mathrm{d}a\,\mathrm{d}s\,\mathrm{d}c,
\qquad
X_1=\partial_s,
\qquad
H:=-iX_1=-i\partial_s.
\tag{7.1}
\label{eq:flowbox}
\end{equation}
\endgroup
Indeed, if \(\Psi\) is a local inverse sheet, \(D\Psi=(DF)^{-1}\) and
\[
\bigl|\det D\Psi\bigr|
=\frac{1}{|\det DF|}=\tfrac12,
\]
so Lebesgue measure pushes forward with factor \(\tfrac12\) (orientation chosen so the factor is positive).

Eliminating \(q_2\) from \(F_2=c\) and \(q_1\) from \((F_0,F_1)=(a,s)\) yields
\begingroup
\renewcommand*{\theHequation}{cubic}
\begin{equation}
A(s;a,c)\,q_0^3+B(s;c)\,q_0+C(c)=0,
\tag{7.2}
\label{eq:cubic}
\end{equation}
\endgroup
with
\[
A=-cs^3+s^2+18acs-27a^2c^2-16a,\quad
B=3cs-4,\quad
C=2c.
\]
If \(A\to 0\) with \(B\neq 0\), Vieta forces some \(|q_0|\to\infty\).

\paragraph{Reconstruction.}
If \(x=q_0\neq 0\) solves \eqref{eq:cubic} and the denominators below are nonzero,
\begingroup
\renewcommand*{\theHequation}{reconstruction}
\begin{equation}
\begin{aligned}
q_1
&=\frac{9acx^2-3csx-3c-sx^2+6x}{x\bigl((3cs-4)x+3c\bigr)},\\
q_2
&=\frac{2x-3x^2q_1-c}{x^3},
\end{aligned}
\tag{8.0}
\label{eq:recon}
\end{equation}
\endgroup
then \(F(x,q_1,q_2)=(a,s,c)\) by polynomial identity.

\section{Escape walls and deficiency indices}
Write
\[
n_+:=\dim_{\mathrm H}\ker(H^*-i),\qquad
n_-:=\dim_{\mathrm H}\ker(H^*+i),
\]
where \(\dim_{\mathrm H}\) denotes Hilbert dimension, i.e.\ orthonormal-basis
cardinality, not algebraic/Hamel dimension.

\subsection{Per-orbit dictionary}
On \(L^2((\ell,r),\mathrm{d}s)\) with \(h=-i\partial_s\) minimal on \(C_c^\infty(\ell,r)\)~\cite{RS2}:
finite upper end \(\Rightarrow\dim_{\mathrm H}\ker(h^*+i)=1\) (model \(e^{s-r}\));
finite lower end \(\Rightarrow\dim_{\mathrm H}\ker(h^*-i)=1\) (model \(e^{-(s-\ell)}\)).

\subsection{Forward wall}
\begin{proposition}[Forward wall]
There exist a nonempty open \(U_+\subset\mathbb{R}^2\) and \(\beta\in C^\infty(U_+)\) such that for every \((a,c)\in U_+\) a selected real inverse branch of \(F\) exists on an \(s\)-interval with finite upper end \(s=\beta(a,c)\) and \(\|q\|\to\infty\) as \(s\to\beta^-\).
\end{proposition}
\begin{proof}
At the base \(p_\star=(a,s,c)=(0,\tfrac12,2)\): \(A(p_\star)=0\), \(A_s:=\partial_s A(p_\star)=-\tfrac12\neq 0\), \(B(p_\star)=-1\neq 0\). By the implicit function theorem there is a neighborhood of \((0,2)\) and \(\beta\in C^\infty\) with \(\beta(0,2)=\tfrac12\) and \(A(\beta(a,c);a,c)=0\).

Set \(\tau=\sqrt{\beta-s}\) and \(q_0=r/\tau\). Define
\[
G_+(a,c,\tau,r)
:=\frac{A(\beta-\tau^2;a,c)}{\tau^2}\,r^3
+B(\beta-\tau^2;c)\,r
+C(c)\,\tau
\quad(\tau\neq 0),
\]
and extend at \(\tau=0\) by
\[
G_+(a,c,0,r)=(-A_s(\beta;a,c))r^3+B(\beta;c)r.
\]
Smoothness of the extension follows from \(A(\beta)=0\) and the Taylor expansion \(A(\beta-\tau^2)=-A_s(\beta)\tau^2+O(\tau^4)\) in \(\tau^2\). At the base, \(G_+(0,2,0,r)=\tfrac12 r^3-r\) has simple roots \(r=\pm\sqrt{2}\) because \(\partial_r G_+=\tfrac32 r^2-1\) and \(\partial_r G_+(\pm\sqrt{2})=-2B=2\neq 0\). The IFT continues \(r_\pm(a,c,\tau)\) on an open set \(U_+\times(0,\tau_0)\). Then \(|q_0|=|r_\pm|/\tau\to\infty\) as \(\tau\downarrow 0\).

For reconstruction, the denominator in \eqref{eq:recon} behaves as \(q_0(Bq_0+3c)\sim B q_0^2\) with \(B\to-1\neq 0\), hence is nonzero for small \(\tau\) after shrinking \(U_+\). The formulas for \(q_1,q_2\) are rational in \((q_0,s,a,c)\) with this nonvanishing denominator; combined with \(q_0\sim r/\tau\) one obtains the spatial escape lower bound
\[
\|q\|\ge |q_0|\ge \frac{c_0}{\sqrt{\beta-s}}
\]
for some \(c_0>0\) on a possibly smaller relatively compact \(W_+\Subset U_+\) and \(s\) near \(\beta\).
Regression certificate (companion release~\cite{FredriksenSoft}):
\filepath{data/anchor/cas_orbit_measure_IFT_A001.json}.
\end{proof}

\subsection{Independent backward-wall geometry (not used in Theorem~F)}
The following calculation is retained as a geometric comparison with the
forward wall. The deficiency proof below does not use it; the \(+i\) space is
obtained from the forward \(-i\) space by the Lean-backed sign involution.
\begin{proposition}[Backward wall]
There exist a nonempty open \(U_-\) and \(\alpha\in C^\infty(U_-)\) such that for \((a,c)\in U_-\) a selected real branch exists for \(s>\alpha(a,c)\) with \(\|q\|\to\infty\) as \(s\downarrow\alpha^+\).
\end{proposition}
\begin{proof}
At \(a=\tfrac1{54}\), \(c=2\),
\[
A\bigl(s;\tfrac1{54},2\bigr)=-\frac{(2s-1)(3s^2-1)}{3},
\]
so \(A(\tfrac12)=0\), \(A_s(\tfrac12)=\tfrac16>0\), \(B(\tfrac12)=-1\neq 0\). By the IFT there is a neighborhood of \((\tfrac1{54},2)\) and \(\alpha\in C^\infty\) with \(\alpha(\tfrac1{54},2)=\tfrac12\) and \(A(\alpha(a,c);a,c)=0\).

Set \(\tau=\sqrt{s-\alpha}\) (so \(s=\alpha+\tau^2\)) and \(q_0=r/\tau\). Define
\[
G_-(a,c,\tau,r)
:=\frac{A(\alpha+\tau^2;a,c)}{\tau^2}\,r^3
+B(\alpha+\tau^2;c)\,r
+C(c)\,\tau
\quad(\tau\neq 0),
\]
and extend at \(\tau=0\) by
\[
G_-(a,c,0,r)=A_s(\alpha;a,c)\,r^3+B(\alpha;c)\,r.
\]
Smoothness of the extension follows from \(A(\alpha)=0\) and \(A(\alpha+\tau^2)=A_s(\alpha)\tau^2+O(\tau^4)\). At the base, \(G_-(\tfrac1{54},2,0,r)=\tfrac16 r^3-r\) has simple roots \(r=\pm\sqrt{6}\) because \(\partial_r G_-=\tfrac12 r^2-1\) and \(\partial_r G_-(\pm\sqrt{6})=-2B=2\neq 0\). The IFT continues \(r_\pm(a,c,\tau)\) on an open set \(U_-\times(0,\tau_0)\), with \(r_\pm\neq 0\) after shrinking. Then \(|q_0|=|r_\pm|/\tau\to\infty\) as \(\tau\downarrow 0\).

For reconstruction, the denominator in \eqref{eq:recon} behaves as \(q_0(Bq_0+3c)\sim B q_0^2\) with \(B\to-1\neq 0\), hence is nonzero for small \(\tau\) after shrinking \(U_-\). Combined with \(q_0\sim r/\tau\) one obtains
\[
\|q\|\ge |q_0|\ge \frac{c_0'}{\sqrt{s-\alpha}}
\]
on a relatively compact \(W_-\Subset U_-\) near the wall.
Regression certificate (companion release~\cite{FredriksenSoft}):
\filepath{data/anchor/cas_backward_incomplete_wall_A001.json}.
\end{proof}

\subsection{Saturated maximal sheets}
\begin{lemma}[Saturated maximal sheets]
Let \(W_+\Subset U_+\) be nonempty, open, and relatively compact. On the selected large-\(q_0\) forward branch, pick \(\varepsilon_0>0\) small and the smooth cross-section
\[
\Sigma_+=\bigl\{q:F(q)=(a,\beta(a,c)-\varepsilon_0,c),\;(a,c)\in W_+\bigr\}.
\]
Let \(\phi_t\) be the maximal flow of \(X_1\). Choose \(\varepsilon_0>0\) small enough that \(\varepsilon_0<\tau_0^2\) for the IFT window of the forward-wall proposition, so that on \(W_+\) the selected large-\(q_0\) branch is a single smooth immersed sheet with nonvanishing reconstruction denominator throughout the open interval \((\beta-\varepsilon_0,\beta)\). Because \(X_1 F_1=1\), \(F_1(\phi_t(q))=F_1(q)+t\) on the existence interval. Define
\[
D_+=\bigl\{(a,c,s):(a,c)\in W_+,\;\ell_+(a,c)<s<\beta(a,c)\bigr\},
\]
where \(\beta\) is the wall value (hence the maximal upper \(s\)-time of this sheet) and \(\ell_+\) is the maximal lower \(s\)-time of the immersed branch through \(\Sigma_+\). Then
\[
\Psi_+:D_+\to\Omega_+\subset\mathbb{R}^3,
\qquad
\Psi_+(a,c,s)=\phi_{s-(\beta-\varepsilon_0)}(q_\Sigma(a,c))
\]
is a \(C^\infty\) diffeomorphism onto an open \(X_1\)-invariant set, \(F\circ\Psi_+=(a,s,c)\), and \(s\to\beta^-\) occurs at spatial infinity. If \(\ell_+\) is finite, maximality forces escape as \(s\downarrow\ell_+\) as well.
\end{lemma}
\begin{proof}
The maximal-flow domain \(\mathcal{D}\subset\mathbb{R}\times\mathbb{R}^3\) of the smooth vector field \(X_1\) is open and the flow \(\phi\) is \(C^\infty\) on \(\mathcal{D}\)~\cite{LeeSmooth}. The set \(D_+\) is (up to the smooth reparameterization \(s=F_1\)) the pullback of \(\mathcal{D}\) through the cross-section \(\Sigma_+\); openness of \(\mathcal{D}\) therefore yields openness of \(D_+\). Along every orbit, \(F_0\) and \(F_2\) are constant and \(F_1\) is the time coordinate. Hence two points of the parameter domain with different \((a,c)\) cannot map to the same orbit, while ODE uniqueness handles two parameterizations on a fixed orbit. Equivalently, the identity \(F\circ\Psi_+=(a,s,c)\) immediately gives injectivity of \(\Psi_+\). Since \(\det DF\neq 0\), \(\Psi_+\) is a local diffeomorphism, hence a \(C^\infty\) diffeomorphism onto its open image \(\Omega_+\). By the \(\varepsilon_0\)-window choice, the selected branch does not degenerate before the wall, so the maximal upper time equals \(\beta\). Finite lower maximal time forces escape from compacts by standard ODE maximality.

No continuity or semicontinuity of the variable endpoint
\(\ell_+(a,c)\) is needed below. In particular, the weak-adjoint argument is
carried out directly on the open variable domain \(D_+\), rather than through
a global assertion about the finite boundary of \(\Omega_+\).
\end{proof}

\paragraph{Lean coverage for the maximal sheet.}
The diffeomorphism and \(X_1\)-invariance formulation in this lemma is the
paper-level maximal-flow packaging. The pinned Lean route formalizes the
ingredients consumed by the deficiency proof, including
\texttt{injOn\_PsiFin3}, \texttt{isOpen\_image\_Psi},
\texttt{measurableEmbedding\_PsiFin3}, the exact coordinate identity, and
\texttt{abs\_det\_fderiv\_PsiFin3}; it does not package a named global
diffeomorphism or \(X_1\)-invariance theorem.

\subsection{Deficiency vectors in
\texorpdfstring{\(\operatorname{Dom}(H^*)\)}{Dom(H*)}}
\begin{lemma}[Deficiency vectors]
For \(\chi\in C_c^\infty(W_+)\), so in particular
\(\operatorname{supp}\chi\Subset W_+\), define on the entire maximal interval
\[
u_-\bigl(\Psi_+(a,c,s)\bigr)
:=\chi(a,c)\,e^{s-\beta(a,c)},
\]
and \(u_-=0\) off \(\Omega_+\). No interior \(s\)-cutoff. Then \(u_-\in L^2(\mathbb{R}^3)\), \(u_-\in\operatorname{Dom}(H^*)\), and \((H^*+i)u_-=0\). The linear map \(\chi\mapsto u_-\) is injective, so \(n_-=\infty\).
\end{lemma}
\begin{proof}
\textit{Square-integrability.} By \eqref{eq:flowbox},
\[
\|u_-\|_2^2
=\tfrac12\int_{W_+}|\chi|^2
\Bigl(\int_{\ell_+}^{\beta}e^{2(s-\beta)}\,\mathrm{d}s\Bigr)
\mathrm{d}a\,\mathrm{d}c
\le\tfrac14\|\chi\|_{L^2(W_+)}^2,
\]
since the inner integral equals \(\tfrac12\bigl(1-e^{2(\ell_+-\beta)}\bigr)\le\tfrac12\).

\textit{Adjoint equation.} Let \(\varphi\in C_c^\infty(\mathbb{R}^3)\) and
put \(K=\operatorname{supp}\chi\Subset W_+\). On \(\Omega_+\),
\(H=-i\partial_s\) in the flow-box chart with density \(\tfrac12\), and
\(\operatorname{div}X_1=0\). We verify the weak identity for the zero extension
directly. This avoids any inference of regularity across
\(\partial\Omega_+\), and it uses no continuity property of \(\ell_+\).

Both terms in the weak-adjoint identity are absolutely integrable. Indeed,
\(X_1\varphi\) is bounded because \(\varphi\) has compact support (the
polynomial growth of \(X_1\) is irrelevant on that compact set), and on
\(D_+\)
\[
\begin{split}
&\bigl|\chi(a,c)e^{s-\beta(a,c)}
  \varphi(\Psi_+(a,c,s))\bigr|\\
&\quad+
\bigl|\chi(a,c)e^{s-\beta(a,c)}
  (X_1\varphi)(\Psi_+(a,c,s))\bigr|
\le
\bigl(\|\varphi\|_\infty+\|X_1\varphi\|_\infty\bigr)
|\chi(a,c)|e^{s-\beta(a,c)} .
\end{split}
\]
The right side is integrable because
\(\int_{\ell_+(a,c)}^{\beta(a,c)}
e^{s-\beta(a,c)}\,\mathrm{d}s\le1\) and \(\chi\) has compact transverse
support.

Change variables on \(\Omega_+\) and apply Fubini; the preceding majorant
proves absolute integrability of both pairings. Now fix \((a,c)\). On every
compact interval
\([r,t]\Subset(\ell_+(a,c),\beta(a,c))\), ordinary one-dimensional
integration by parts applies. Let \(t\uparrow\beta(a,c)\) and
\(r\downarrow\ell_+(a,c)\) separately. These are limits on one fixed fiber:
the two improper integrals converge absolutely by the same exponential
estimate. At the upper endpoint the sheet escapes every compact set, so
\(\varphi\) vanishes eventually. At the lower endpoint, either \(\ell_+\) is
finite and maximality again forces spatial escape, or \(\ell_+=-\infty\) and
\(e^{s-\beta}\to0\). Hence both endpoint terms vanish, and the desired
one-dimensional weak identity holds on that fiber.

We then integrate this already established fiberwise equality over \((a,c)\).
Fubini applies to both sides by the global majorant above; there is no
transverse limiting process, measurable exhaustion, or endpoint-selection
argument. Consequently
\[
\langle u_-,H\varphi\rangle=\langle -iu_-,\varphi\rangle,
\]
i.e.\ \(u_-\in\operatorname{Dom}(H^*)\) and \((H^*+i)u_-=0\).

\textit{Injectivity with a uniform bound.} The collar used to define
\(\Sigma_+\) lies in every selected maximal interval, so
\(\ell_+(a,c)<\beta(a,c)-\varepsilon_0\). Therefore
\[
\|u_-\|_2^2
\ge \frac14\bigl(1-e^{-2\varepsilon_0}\bigr)
\|\chi\|_{L^2(W_+)}^2 .
\]
Thus \(\chi\mapsto u_-\) is bounded below, hence injective, and embeds the
infinite-dimensional space \(C_c^\infty(W_+)\) into \(\ker(H^*+i)\).
\end{proof}

\paragraph{Lean coverage for the cutoff family.}
The displayed bounded-below map and its uniform quantitative constant are
paper-level. The pinned Lean proof reaches the same required
non-finite-dimensional conclusion through
\texttt{exists\_finite\_weakAdjointEigenvector\_negI\_family}: for every
finite cardinality it constructs linearly independent deficiency vectors
from disjoint compactly supported transverse cutoffs.

\subsection{Sign involution and the opposite deficiency space}
\begin{lemma}[Sign transport]
Let
\(\sigma(q_0,q_1,q_2)=(-q_0,-q_1,q_2)\) and
\((U_\sigma f)(q)=f(\sigma(q))\). Then \(U_\sigma\) is a complex-linear
unitary involution of \(L^2(\mathbb{R}^3)\), preserves
\(C_c^\infty(\mathbb{R}^3)\), and
\[
H U_\sigma=-U_\sigma H
\quad\hbox{on }C_c^\infty(\mathbb{R}^3).
\]
Consequently \(H^*U_\sigma=-U_\sigma H^*\), and \(U_\sigma\) restricts to a
complex-linear isometric bijection
\[
\ker(H^*+i)\;\xrightarrow{\;\cong\;}\;\ker(H^*-i).
\]
\end{lemma}
\begin{proof}
The linear involution \(\sigma\) preserves Lebesgue measure, so its pullback is
unitary and preserves the test-function core. Direct substitution in the
explicit polynomial field gives
\[
X_1(f\circ\sigma)=-(X_1f)\circ\sigma,
\]
which is the displayed core anticommutation for \(H=-iX_1\). For
\(v\in\operatorname{Dom}(H^*)\) and a core test function, unitarity and core
anticommutation transfer the defining adjoint pairing and give
\(U_\sigma v\in\operatorname{Dom}(H^*)\) with
\(H^*U_\sigma v=-U_\sigma H^*v\). Since \(U_\sigma^2=I\), this transports the
\(-i\) eigenspace isometrically and bijectively to the \(+i\) eigenspace. This
is the route formalized for the specific canonical minimal operator
\texttt{H\_X1\_min} in the companion Lean artifact.
\end{proof}

\begin{remark}[Withdrawn device]
Interior indicators \(\mathbf{1}_{(\beta-\delta,\beta)}(s)\) create \(\delta\)-masses and exit \(\operatorname{Dom}(H^*)\). That construction is withdrawn.
\end{remark}

\begin{theorem}[F]
The Hilbert deficiency indices satisfy
\[
n_+=n_-=\aleph_0,
\]
equivalently \((n_+,n_-)=(\infty,\infty)\) in conventional notation.
\end{theorem}
\begin{proof}
The forward-wall deficiency lemma makes \(\ker(H^*+i)\)
infinite-dimensional. Sign transport gives an isometric bijection with
\(\ker(H^*-i)\), so the opposite space is infinite-dimensional as well.
Both are closed subspaces of the separable Hilbert space
\(L^2(\mathbb{R}^3)\); their Hilbert dimensions (orthonormal-basis cardinals)
are therefore exactly \(\aleph_0\). This is a Hilbert-dimension statement, not
an assertion that the algebraic Hamel rank is \(\aleph_0\).
\end{proof}

\begin{theorem}[E]
\(H=-iX_1\) is not essentially self-adjoint on \(C_c^\infty(\mathbb{R}^3)\).
\end{theorem}
\begin{proof}
\(n_\pm\ge 1\) by Theorem~F.
\end{proof}

\begin{corollary}
Let \(\overline H\) denote the closure of \(H\). Von Neumann's theorem applies
to the closed symmetric operator \(\overline H\). Since
\(H^*=(\overline H)^*\), its deficiency spaces are those computed above;
moreover, a self-adjoint operator extends \(H\) if and only if it extends
\(\overline H\). Fixing one deficiency-space isometry and multiplying it by
the unit phases \(\lambda\in\mathbb C\), \(|\lambda|=1\), gives distinct
extensions under the von Neumann bijection. Lean proves the injectivity of
this phase parameterization; the classical fact that the complex unit circle
has cardinality continuum then gives a continuum-sized family of
self-adjoint extensions. Every
such extension agrees with the same minimal operator on
\(C_c^\infty(\mathbb{R}^3)\); the core algebraic Poisson/CCR relations
recorded in Theorems~B--C therefore do not, by themselves, distinguish a
unique extension.
\end{corollary}

\section{Discussion}
The example separates algebraic and analytic layers. The Poisson cotangent lift and dual-field/CCR identities are clean, and \(\psi\) is a well-defined polynomial Weyl endomorphism. Nonetheless \(X_1\) is incomplete and \(H=-iX_1\) has Hilbert deficiency indices \((\aleph_0,\aleph_0)\). Algebraic CCR/Poisson data therefore do not select a unique self-adjoint dual momentum for this generator.

We do \emph{not} claim here exact deficiency indices for \(H_0\) or \(H_2\), nor a theorem classifying all strong CCR or dual-translation packages; those topics appear only in companion repository notes and may be incomplete relative to theorem grade.

The collision is used in Theorem~A and in the noninjectivity clause of
Theorem~B, but it is not an input to Theorems~D--F. Those analytic results use
the constant nonzero Jacobian (to define the dual field and its invariant
measure) and the explicit incompleteness/escape geometry of \(X_1\). Thus the
logical claim is not ``counterexample implies quantization ambiguity.''
Classical incomplete flows, half-line deficiency models, and von Neumann
extension theory are standard~\cite{RS2}. The contribution is their explicit,
Lean-backed realization for this dual transport operator. We make no new
general theorem equating surjectivity of a real Keller map with completeness
of all its dual fields.

\section{Non-claims}
No unitary gate, channel, CP map, or computational advantage. No preferred
physical extension. No claim that all quantization routes fail. Discovery of
\(F\) is not claimed. No rank-three Dixmier counterexample is claimed. Lean
covers the bounded canonical-core Theorems E--F, the exact Hilbert-basis
cardinal result, the specific \texttt{H\_X1\_min} von Neumann classification,
and the injective unit-phase family of distinct extension witnesses at the
synchronized freeze identified above. The algebraic index
\texttt{hamelDeficiency}\allowbreak\texttt{Rank}, defined using
\texttt{Module.rank}, has lower bounds only.
The independent backward-wall construction, exact cardinality of the full
extension type, arbitrary-operator theorem, strong CCR, preferred extension,
and algebraic or physical selection remain outside the Lean corollary. No
theorem-level dual-flow strong CCR obstruction is claimed in this note.
In particular, the artifact constructs the displayed Weyl endomorphism
\(\psi\) but does not establish that \(\psi\) is noninvertible.

\footnotesize
\sloppy
\raggedright
\begin{thebibliography}{99}
\bibitem{RS2}
M.~Reed and B.~Simon, \emph{Methods of Modern Mathematical Physics II: Fourier Analysis, Self-Adjointness}, Academic Press, 1975.

\bibitem{DixmierWeyl}
J.~Dixmier, \emph{Sur les alg\`ebres de Weyl}, Bull.\ Soc.\ Math.\ France \textbf{96} (1968), 209--242.

\bibitem{Tsuchimoto2005}
Y.~Tsuchimoto, \emph{Endomorphisms of Weyl algebra and
\(p\)-curvatures}, Osaka J.\ Math.\ \textbf{42} (2005), 435--452,
\url{https://doi.org/10.18910/7472}.

\bibitem{BKK2007}
A.~Belov--Kanel and M.~Kontsevich, \emph{The Jacobian conjecture is
stably equivalent to the Dixmier conjecture}, Moscow Math.\ J.\
\textbf{7} (2007), 209--218,
\url{https://arxiv.org/abs/math/0512171}.

\bibitem{AdjamagboEssen2007}
P.~K.~Adjamagbo and A.~van den Essen, \emph{A proof of the equivalence of
the Dixmier, Jacobian and Poisson conjectures}, Acta Math.\ Vietnam.\
\textbf{32} (2007), 205--214,
\url{https://math.ac.vn/uploads/files/0702205.pdf}.

\bibitem{LeeSmooth}
J.~M.~Lee, \emph{Introduction to Smooth Manifolds}, 2nd ed., Springer, 2013.

\bibitem{BassCW}
H.~Bass, E.~H.~Connell, and D.~Wright, The Jacobian conjecture: reduction of degree and formal expansion of the inverse, Bull.\ Amer.\ Math.\ Soc.\ \textbf{7} (1982), 287--330.

\bibitem{Alpoge2026}
L.~Alp\"oge, informal public announcement of a Keller map with constant
Jacobian and three-point collision (credits Akhil Mathew for prompting the
question and Claude Fable for work leading to the map), July 2026,
\url{https://x.com/__alpoge__/status/2079028340955197566}.

\bibitem{UlamJacobian}
Ulam, \emph{A Counterexample to the Jacobian Conjecture}, preprint, July 2026,
\url{https://www.ulam.ai/research/jacobian.pdf}.

\bibitem{AFPJacobian}
A.~Freitas Ramos, D.~Barros Hulak, and R.~J.~G.~Barretto de Queiroz,
\emph{Formal Verification of an Explicit Counterexample to the Jacobian
Conjecture}, Archive of Formal Proofs, 20 July 2026,
\url{https://isa-afp.org/entries/Jacobian_Counterexample.html}.

\bibitem{Shaska2026}
T.~Shaska, \emph{Graded Keller maps and the Jacobian Conjecture},
arXiv:2607.20210, 2026,
\url{https://arxiv.org/abs/2607.20210}.

\bibitem{JelonekLason2018}
Z.~Jelonek and M.~Laso\'n, \emph{Quantitative properties of the
non-properness set of a polynomial map}, Manuscripta Math.\
\textbf{156} (2018), 383--397,
\url{https://doi.org/10.1007/s00229-017-0965-0}.

\bibitem{SpeyerJC2026}
D.~Speyer, \emph{The new counterexample to the Jacobian conjecture}, Secret Blogging Seminar, 20 July 2026,
\url{https://sbseminar.wordpress.com/2026/07/20/the-new-counterexample-to-the-jacobian-conjecture/}.

\bibitem{FredriksenSoft}
D.~E.~Fredriksen, EXOTIC-CCR A001 companion code and CAS certificates,
GitHub Quantyra/jacobian-weyl-quantum-phase-space, proposed paper-package
candidate v0.3.9-referee-revision (untagged);
concept DOI \url{https://doi.org/10.5281/zenodo.21474351} (project-level).

\bibitem{FredriksenLean}
D.~E.~Fredriksen, exotic-ccr-lean, GitHub Quantyra/exotic-ccr-lean,
released tag v0.1.8-theorem-e, tag target be4f330; reviewed bounded
Theorem~E source anchor a6bb091,
\url{https://github.com/Quantyra/exotic-ccr-lean/releases/tag/v0.1.8-theorem-e}
(Lean~4 formalization of the seed identities, dual-matrix algebra, Piola row
divergence, full dual-field commutation, the polynomial Weyl endomorphism
\(\psi\), and incompleteness of the dual field \(X_1\)).
\bibitem{FredriksenLeanF}
Theorem F supplement: synchronized, untagged, unreleased Lean publication
freeze fbcdd0345d2f2540cd537204be2178ae07e18a5e (historical theorem-source
root ff50f4a2a312591c2e5b26e71eb390ade9164b34);
\path{ExoticCCR.theoremF},
\path{aleph0_le_hamelDeficiencyRank_X1}, and
\path{hilbertDeficiencyIndex_X1_eq_aleph0},
\texttt{theoremFVonNeumannClassification},
\texttt{theoremFUnitPhaseExtension}, and
\texttt{theoremFUnitPhaseExtension\_injective}. Both standard
adjoint eigenspaces are not finite-dimensional, and
\(\operatorname{Cardinal.aleph0}\) is a lower bound for each
\texttt{hamelDeficiencyRank}. The separate Hilbert-basis
cardinal is exactly \(\operatorname{aleph0}\) for both closed eigenspaces;
\texttt{hamelDeficiencyRank} is the algebraic (Hamel) rank and its exact
value is not claimed. For the specific canonical minimal operator, all
self-adjoint-extension witnesses are classified bijectively by complex-linear
isometric equivalences from the $+i$ to the $-i$ adjoint eigenspace, and
unitary complex phases inject into distinct witnesses. Together with the
classical cardinality of the unit circle, this yields a continuum-sized lower
family; the cardinal identification is not asserted as a Lean declaration.
No exact cardinality of the full extension type,
preferred extension, arbitrary-operator theorem, backward-wall construction,
or algebraic/physical selection theorem is claimed.
\end{thebibliography}
\end{document}
